\documentclass{conjura-conjecture}

\newcommand{\Fun}{\mathsf{Fun}}
\newcommand{\Adv}{\mathbf{Adv}}
\newcommand{\Gm}{\mathsf{Game}}
\newcommand{\Pred}{\mathrm{Pred}}
\newcommand{\ExtP}{\mathrm{Ext}\text{-}\mathrm{pub}}
\newcommand{\pred}{\mathrm{pred}}
\newcommand{\extpub}{\textup{ext-pub}}
\newcommand{\sd}{\mathit{sd}}
\renewcommand{\sample}{\stackrel{{\scriptscriptstyle\$}}{\leftarrow}}
\newcommand{\SD}{\mathrm{SD}}
\newcommand{\calK}{\mathcal{K}}
\newcommand{\calD}{\mathcal{D}}
\newcommand{\calR}{\mathcal{R}}
\newcommand{\ret}{\textbf{return}\ }

\runninghead{LEFTOVER HASH LEMMA CONJECTURE, PUBLIC SEED}

\begin{document}

\cjkicker{CONJURA \textperiodcentered{} OPEN PROBLEM}
\cjtitle{Randomness Extraction from Unpredictable Random-Oracle Sources}
\cjsubtitle{A Leftover-Hash-Lemma Conjecture, Public Seed}
\cjstatus{Statement: human-written, formalized in Lean (see \texttt{lean/Statement.lean}).}
\cjcategory{Information-Theoretic Cryptography.}

\vspace{0.4em}

\noindent\textbf{Setting.}
Fix nonempty finite sets $\calK$ (seeds), $\calD$ (inputs), $\calR$ (outputs), and write
$K := |\calK|$, $D := |\calD|$, $R := |\calR|$. Let $\Fun(\calK \times \calD, \calR)$ be the set of
all functions $\calK \times \calD \to \calR$, let $\SD(\cdot,\cdot)$ denote statistical distance and
$U_{\calR}$ the uniform distribution on $\calR$. The source $S$, the predictor $P$ and the
distinguisher $\mathsf{D}$ of Figure~\ref{fig:games} are computationally unbounded and receive the
\emph{entire} function table of $H$ as an explicit input. The seed $\sd$ is sampled only after $S$
has terminated, and is given to $\mathsf{D}$: that is what makes this the \emph{public-seed} game.

\begin{figure}[ht]
\centering
\fbox{\begin{minipage}[t]{0.85\linewidth}
\underline{$\Gm\ \Pred^{S}_{P}$}\\[2pt]
$H \sample \Fun(\calK \times \calD, \calR)$;\quad
$(x,z) \sample S(H)$;\quad
$x' \sample P(H,z)$\\
$\ret (x = x')$
\end{minipage}}

\smallskip
\fbox{\begin{minipage}[t]{0.85\linewidth}
\underline{$\Gm\ \ExtP^{S}_{\mathsf{D}}$ (public seed)}\\[2pt]
$H \sample \Fun(\calK \times \calD, \calR)$;\quad
$(x,z) \sample S(H)$;\quad $\sd \sample \calK$\\
$y_0 \leftarrow H(\sd,x)$;\quad $y_1 \sample \calR$;\quad
$b \sample \{0,1\}$\\
$b' \sample \mathsf{D}(H, \sd, y_b, z)$;\quad
$\ret (b = b')$
\end{minipage}}
\caption{Prediction game (top) and the public-seed extraction game (bottom). The first argument of
$S$, $P$ and $\mathsf{D}$ is the \emph{full function table} of $H$; the distinguisher additionally
receives the seed $\sd$.}
\label{fig:games}
\end{figure}

\begin{definition}\label{def:games}
Set $\Adv^{\pred}_{\calK,\calD,\calR,S}(P) := \Pr[\Pred^{S}_{P} \Rightarrow 1]$ and
\[
  \Adv^{\extpub}_{\calK,\calD,\calR}(S,\mathsf{D}) := 2\Pr[\ExtP^{S}_{\mathsf{D}} \Rightarrow 1] - 1 .
\]
A source $S$ is \emph{$\epsilon$-unpredictable} if
$\Adv^{\pred}_{\calK,\calD,\calR,S}(P) \le \epsilon$ for every unbounded predictor $P$.
\end{definition}

\begin{conjecture}[Public seed]\label{conj:pub}
There is a universal constant $c > 0$, independent of $\calK$, $\calD$, $\calR$ and $\epsilon$, such
that for all nonempty finite $\calK, \calD, \calR$, all $\epsilon \in (0,1]$, every
$\epsilon$-unpredictable source $S$ and every unbounded distinguisher $\mathsf{D}$,
\[
  \Adv^{\extpub}_{\calK,\calD,\calR}(S,\mathsf{D})
  \;\le\; \delta_{\mathrm{pub}}(\epsilon, K, D, R)
  \;:=\; c \cdot \Bigl( \sqrt{\epsilon R} + \sqrt{\frac{\log_2 D}{K}} \,\Bigr) .
\]
\end{conjecture}

\begin{remark}[What the leftover hash lemma gives]\label{rem:lhl}
Condition on the table and the auxiliary output, $(H,Z) = (h,z)$, and write
$p_x := \Pr[X = x \mid h,z]$. Because $\sd$ is drawn only after $S$ has terminated it is uniform on
$\calK$ and independent of $X$ given $(h,z)$, so the distinguisher's task is to tell
$(\sd, h(\sd,X))$ from uniform on $\calK \times \calR$: a strong-extractor question about the
\emph{fixed} family $\{h(k,\cdot)\}_{k \in \calK}$, to which the leftover hash lemma applies in its
almost-universal form. Writing $\mathrm{CP} := \sum_x p_x^2$ for the source's collision probability
and $\bar\beta := \tfrac1K \sum_{k \in \calK} \Pr[h(k,X) = h(k,X') \mid X \neq X']$ for the family's
mean collision probability on that source, it bounds the conditional advantage by
\[
  \tfrac12\sqrt{\,R\cdot\mathrm{CP} \;+\; \gamma\,},
  \qquad \gamma \;:=\; R\bar\beta - 1 ,
\]
where $\gamma \le 0$ exactly when the family is universal on the source's support. Averaging over
$(H,Z)$ preserves the shape, $\sqrt{\cdot}$ being concave, and $\mathrm{CP} \le \max_x p_x$ turns
the first summand into $\tfrac12\sqrt{\epsilon R}$.

Two things follow, and between them they say what this conjecture asks for that the lemma does not
supply. On the first summand the lemma already delivers the conjectured shape, with $c = \tfrac12$.
On the second it cannot. The family $\{h(k,\cdot)\}$ is not universal and cannot be made so: a
universal family on $D$ inputs needs a seed of $\Omega(\log D)$ bits, that is $K \ge D^{\Omega(1)}$,
whereas the conjecture is of interest at $K \approx \log_2 D$ (Vadhan, \emph{Pseudorandomness},
the remark following Theorem~6.18). So $\gamma$ must be certified for a random table, uniformly over
the supports $S$ may pick \emph{after} seeing $H$, and whatever certification is used it enters
under the square root. Since $\bar\beta$ is a collision frequency read off the $Kt$ table entries
above a support of size $t$, and the estimate must hold uniformly over $\binom{D}{t}$ of them, it
cannot be pinned closer than $\sqrt{\log D / K}$; so the lemma yields
\[
  \Adv^{\extpub}_{\calK,\calD,\calR}(S,\mathsf{D})
  \;\lesssim\; \tfrac12\sqrt{\epsilon R}
  \;+\; O\!\left(\Bigl(\frac{\log_2 D}{K}\Bigr)^{\!1/4}\right) ,
\]
a fourth root where $\delta_{\mathrm{pub}}$ asks for a square root. In seed length that is
$\log_2 K = \log_2\log_2 D + 4\log_2(1/\delta)$ against the conjectured
$\log_2\log_2 D + 2\log_2(1/\delta)$, the latter being optimal for any extractor
(Radhakrishnan and Ta-Shma, \emph{SIAM J.\ Discrete Math.}\ 13(1), pp.~2--24, 2000).

The conjecture is therefore not a corollary of the leftover hash lemma, and its whole content is the
second summand: the price of letting the source choose its support after seeing the entire table.
The solution note in this folder closes that gap by never passing through a collision probability,
bounding the $L_1$ distance directly instead.
\end{remark}

\bigskip
\noindent\textbf{The secret-seed companion.}
This note states one conjecture, the one the accompanying page states. There is a companion
statement in the same setting, differing only in that $\mathsf{D}$ is called as
$\mathsf{D}(H, y_b, z)$ and never receives $\sd$; it conjectures the bound
$c\sqrt{(\epsilon R + \log_2 D)/K}$, in which the seed's factor $1/K$ attaches to both summands.
It is a genuinely different claim, not a restatement: the solution note in this folder proves the
conjecture above and separately shows that transposing the secret-seed expression verbatim into the
public-seed game gives a \emph{false} statement. The two were originally posed together in a single
write-up and were split so that each note states exactly the conjecture its page states.

\end{document}
